\documentclass{article}

\usepackage{amsmath,amssymb}
\usepackage{xurl}
\usepackage[hidelinks]{hyperref}
\setlength{\emergencystretch}{2em}

% Shared commands for notes in docs/paper-gaps/.

\DeclareUnicodeCharacter{2080}{\ifmmode{}_{0}\else\textsubscript{0}\fi}
\DeclareUnicodeCharacter{2081}{\ifmmode{}_{1}\else\textsubscript{1}\fi}
\DeclareUnicodeCharacter{2082}{\ifmmode{}_{2}\else\textsubscript{2}\fi}
\DeclareUnicodeCharacter{2099}{\ifmmode{}_{n}\else\textsubscript{n}\fi}

\newcommand{\Lean}{\textsc{Lean}}
\ExplSyntaxOn
\NewDocumentCommand{\leanid}{m}{
  \ifmmode
    \textsf{\detokenize{#1}}
  \else
    \group_begin:
    \urlstyle{sf}
    \tl_set:Nn \l_tmpa_tl {#1}
    \tl_replace_all:Nnn \l_tmpa_tl {\_} {_}
    \exp_args:NV \path \l_tmpa_tl
    \group_end:
  \fi
}
\ExplSyntaxOff
\providecommand{\tr}{\operatorname{tr}}
\newcommand{\tnleanIssueBase}{https://github.com/LionSR/TNLean/issues/}
\newcommand{\tnleanPrBase}{https://github.com/LionSR/TNLean/pull/}
\newcommand{\tnleanDocBase}{https://sirui-lu.com/TNLean/docs/}
\newcommand{\ghissue}[1]{\href{\tnleanIssueBase#1}{GitHub issue \##1}}
\newcommand{\ghpr}[1]{\href{\tnleanPrBase#1}{PR \##1}}
\newcommand{\doclink}[1]{\href{\tnleanDocBase#1.html}{\path{#1}}}

% Dirac notation and the identity operator, as in blueprint/src/macros/common.tex.
% \providecommand keeps note-local definitions authoritative.
\usepackage{dsfont}
\providecommand{\ket}[1]{|#1\rangle}
\providecommand{\bra}[1]{\langle #1|}
\providecommand{\braket}[2]{\langle #1 | #2 \rangle}
\providecommand{\ketbra}[2]{|#1\rangle\!\langle #2|}
\providecommand{\Id}{\mathds{1}}
\providecommand{\one}{\mathds{1}}
\providecommand{\C}{\mathbb{C}}
\providecommand{\MN}[1]{M_{#1}(\C)}
\providecommand{\MD}{M_D(\C)}

% External referencing for paper sources.  The build helper writes sanitized
% auxiliary files under build/paper-aux/ and excludes duplicated source labels.
\usepackage{xr}

% Import a generated paper auxiliary file with a caller-supplied label prefix.
% The candidate paths support builds from docs/paper-gaps/, the repository root,
% and build/paper-aux/ itself.
\newcommand{\importpaperaux}[2]{%
  \IfFileExists{../../build/paper-aux/#2.aux}{%
    \externaldocument[#1]{../../build/paper-aux/#2}%
  }{%
    \IfFileExists{build/paper-aux/#2.aux}{%
      \externaldocument[#1]{build/paper-aux/#2}%
    }{%
      \IfFileExists{#2.aux}{%
        \externaldocument[#1]{#2}%
      }{}%
    }%
  }%
}

% General external reference macros
\newcommand{\paperref}[2]{\ref{#1-#2}}
\newcommand{\papereqref}[2]{(\ref{#1-#2})}

% CPSV16 source (1606.00608 / MPDO-22-12-17-2.tex) external document and shortcuts
\importpaperaux{cpsv16-}{1606.00608/MPDO-22-12-17-2-tnlean}
\newcommand{\cpsvref}[1]{\paperref{cpsv16}{#1}}
\newcommand{\cpsveqref}[1]{\papereqref{cpsv16}{#1}}

% Verdict marker: \gapnote{<kind>}{<status>}. Kind is the policy
% classification (clarification, local-correction, scope-restriction,
% unfaithful, false-source, open-gap); status is open, wip (elimination
% underway), resolved, or historical. Machine-read by the site index;
% prints a verdict line.
\newcommand{\gapnote}[2]{\par\noindent\textbf{Verdict:} #1 (#2).\par}


\title{CPSV16 Lemma~\texttt{equalMPS} Gauge-Phase Components and Resolution}
\author{}
\date{2026-05-10}

\begin{document}
\maketitle
\gapnote{unfaithful}{resolved}

\section{The source statement}

Cirac--P\'erez-Garc\'ia--Schuch--Verstraete state the following structural
result for overlaps of normal matrix product vectors
\cite[Lemma~\texttt{equalMPS}, line~1081]{Cirac2016MPDO_arXiv}:

\begin{quote}
    Given two NMPV $V_{a,b}$ generated by two NT with corresponding
    $D_\alpha\times D_\alpha$ matrices $A_\alpha^i$ ($\alpha=a,b$):
    \begin{enumerate}
        \item $\lim_{N\to\infty}\langle V_\alpha\,|\,V_\alpha\rangle = 1$,
        \item $\lim_{N\to\infty}|\langle V_b\,|\,V_a\rangle| = 0$ or $1$.
    \end{enumerate}
    In the latter case, $D_a=D_b$ and there exists a non-singular matrix $X$
    and a phase $\phi$ such that $A_b^i = e^{i\phi} X A_a^i X^{-1}$.
\end{quote}

The source assumes only that the two tensors are normal. Its conclusion has
three parts: the self-overlaps converge to one, the absolute mixed overlap
converges to zero or one, and the unit-overlap case yields equality of bond
dimensions together with gauge-phase equivalence. It does not assume that the
matrix product vector families are proportional.

For tensors $A_a=(A_a^i)_{i=1}^d$ and $A_b=(A_b^i)_{i=1}^d$, define the mixed
transfer matrix by
\begin{align}
    \mathbb E_{ab}
        &= \sum_{i=1}^{d} A_a^i\otimes\overline{A_b^i}.
    \label{eq:cpsv16_equalMPS_gauge_phase_gap_cross_transfer}
\end{align}
The source proof analyzes the eigenvalue relation
\begin{align}
    \sum_{i=1}^{d}(A_b^i)^\dagger X A_a^i
        &= \lambda X
    \label{eq:cpsv16_equalMPS_gauge_phase_gap_eigenvalue_relation}
\end{align}
together with trace-preserving normalization and the strictly positive
spectral fixed point $\Lambda_a$
\cite[lines~1093--1117]{Cirac2016MPDO_arXiv}. Cauchy--Schwarz applied to the
resulting trace inequality gives either $|\lambda|<1$ or $|\lambda|=1$. In the
second case, the equality condition gives
\begin{align}
    A_b^iX
        &= \alpha XA_a^i,
    \label{eq:cpsv16_equalMPS_gauge_phase_gap_intertwining_relation}\\
    |\alpha|
        &= 1,
    \label{eq:cpsv16_equalMPS_gauge_phase_gap_phase_modulus}
\end{align}
and $X$ is an isometry. A dimension argument then gives the structural
conclusion.

\section{The formalized same-dimension component}

The formalization separates the same-bond-dimension part of the unit-overlap
case into two steps. First, asymptotically unit-modulus overlap forces the
spectral radius of the mixed transfer matrix to be at least one.\footnote{The
formal theorem is
\leanid{MPSTensor.mixedTransferSpectralRadius_ge_one_of_mpvOverlap_norm_tendsto_one}
in \path{TNLean/MPS/SharedInfra/GaugePhase.lean}. It formalizes a
step in the proof of Lemma~\texttt{equalMPS}
\cite[lines~1093--1117]{Cirac2016MPDO_arXiv}.}
This statement is the contrapositive of the formalized decay theorem for mixed
transfer spectral radius strictly less than one.

Second, the modulus-one spectral theorem applies the Cauchy--Schwarz equality
argument with the positive fixed point and obtains gauge-phase
equivalence.\footnote{The formal theorem is
\leanid{MPSTensor.modulus_one_eigenvalue_implies_gauge_of_irreducible_TP}
in \path{TNLean/Spectral/TransferOperatorGapNT.lean}.}
Combining these steps gives the same-dimensional gauge-phase theorem with the
source overlap hypothesis and without an MPV proportionality
hypothesis.\footnote{The formal theorem is
\leanid{MPSTensor.gaugePhaseEquiv_of_overlap_norm_tendsto_one_of_irreducible_TP}
in \path{TNLean/MPS/SharedInfra/GaugePhase.lean}.}

This same-dimensional theorem modifies the normal-tensor hypotheses in a
controlled way. In the source, a normal tensor has no nontrivial invariant
projector, and its associated completely positive map has a unique peripheral
eigenvalue, equal to its spectral radius and normalized to one
\cite{Cirac2016MPDO_arXiv}. The formal spectral-radius and Cauchy--Schwarz
argument does not require peripheral simplicity; it assumes irreducibility and
trace-preserving normalization explicitly. It therefore generalizes the
same-dimensional gauge step, but it does not prove the rectangular
bond-dimension conclusion.

The proportionality-based theorem in the same formal file is a valid corollary
under a stronger hypothesis. It is not a formalization of
Lemma~\texttt{equalMPS}, because the source does not assume proportional
matrix product vector families.

\section{The formalized rectangular component}

Lemma~\texttt{equalMPS} also concludes that $D_a=D_b$ in the unit-overlap case
\cite[Lemma~\texttt{equalMPS}]{Cirac2016MPDO_arXiv}. This conclusion is
rectangular: the two normal tensors may have different bond dimensions before
the conclusion is established. The source proof constructs an isometry between
their bond spaces and uses normality of the larger tensor to exclude a proper
missing subspace, thereby forcing equality of dimensions
\cite[lines~1115--1117]{Cirac2016MPDO_arXiv}.

The rectangular overlap-decay theorem states that unequal bond dimensions
force the mixed overlap to converge to zero. Its contrapositive gives
\begin{align}
    \left|\braket{V^{(N)}(B)}{V^{(N)}(A)}\right|
        \xrightarrow[N\to\infty]{}1
    \quad &\Longrightarrow\quad
    D_a=D_b.
    \label{eq:cpsv16_equalMPS_gauge_phase_gap_dimension_recovery}
\end{align}
The decay theorem is
\leanid{MPSTensor.mpvOverlap_tendsto_zero_of_dim_ne_of_irreducible_TP}
in \path{TNLean/Spectral/TransferOperatorGapNT.lean}. Its contrapositive is
recorded as
\leanid{MPSTensor.dim_eq_of_mpvOverlap_norm_tendsto_one_of_irreducible_TP}
in \path{TNLean/MPS/SharedInfra/GaugePhase.lean}. These results use
only positive-length asymptotic overlaps and do not depend on the empty-word
coefficient.

The rectangular conclusion is necessary in the proof of
Theorem~\texttt{thm1}
\cite[statement lines~1167--1170; proof line~1182]{Cirac2016MPDO_arXiv}.
For source blocks $A_j$ and $B_k$, the relevant quantities are
\begin{align}
    \braket{V^{(N)}(B_k)}{V^{(N)}(A_j)}.
    \label{eq:cpsv16_equalMPS_gauge_phase_gap_block_overlap}
\end{align}
The source applies Lemma~\texttt{equalMPS} before matching bond dimensions
have been established. Consequently, the same-dimensional theorem supplies
the gauge step only after the rectangular theorem has supplied
Eq.~\ref{eq:cpsv16_equalMPS_gauge_phase_gap_dimension_recovery}.

\section{Verdict}

The formal development now proves the same-bond-dimension gauge component
without the source-absent MPV proportionality hypothesis. It obtains the
rectangular conclusion $D_a=D_b$ from the non-decay contrapositive of the
rectangular overlap-decay theorem.

The theorem
\leanid{MPSTensor.IsNormalTensor.overlap_dichotomy} in
\path{TNLean/MPS/Overlap/NormalTensorDichotomy.lean} combines these components
with normalized self-overlap and proves the full normal-tensor dichotomy and
structural conclusion of Lemma~\texttt{equalMPS}
\cite{Cirac2016MPDO_arXiv}. Its corollary
\leanid{MPSTensor.IsNormalTensor.mpv_phase_alternative} proves the exact phase
relation at every chain length, as stated in
Corollary~\texttt{eqV} of Ref.~\cite{Cirac2016MPDO_arXiv}. The gap is
resolved; the decomposition above records the mathematical structure of the
formal proof.

\bibliographystyle{alpha}
\bibliography{references}

\end{document}
